Source : calculs de l'auteur (fiscalite/is/calculs\_note\_is.py, mécanique de fiscalite/is/simulation\_is.py).\par Note : taux effectif marginal d'une machine qui rapporte tout juste ce qu'exigent ses financeurs (méthode de Devereux et Griffith, 2003) : financement par fonds propres, rendement réel exigé de 5 \%, inflation de 2 \%, usure de 14 \% par an, amortissement linéaire sur sept ans aujourd'hui. Ces conventions ne sont pas celles de l'OCDE (9,8 \% pour les équipements en 2025) : seuls les écarts entre barres comptent. Coûts en 2027 : la baisse du taux normal, taux réduit gardé, coûte 2,66 Md\euro{} par point avant toute réaction des bénéfices déclarés ; la déduction immédiate de tous les équipements coûte 1,7 à 2,6 Md\euro{} par année d'achats en valeur actualisée au taux des entreprises (5 \%). Le taux de 24 \% est la baisse qui coûte autant.
